% MT002 / explanation — canonical editable content.
题中同一个模型有两种表现：做见过的训练题很准，换成没有参与训练的测试题却不准。这通常提示过拟合（overfitting）：模型连训练数据里的偶然细节也学了进去。比如按几个客户的特殊习惯制定规则，换一批客户就失效。这里讨论的是可能有用的改进，不要求某项操作保证提分。



\begin{itemize}

\item \textbf{A不选。} 随机删除训练样本通常使可用信息更少，小样本中的偶然规律反而更突出；它没有针对题中的问题。

\item \textbf{B不选。} 更复杂的模型能拟合更多细节，但目前训练已经很好，继续增强拟合能力可能扩大训练与测试的差距。

\item \textbf{C选。} 正则化（regularization）是在训练目标中给过大的权重或过高的复杂度增加代价。模型必须为额外的细节付出代价，因此更可能保留能推广的规律。

\item \textbf{D选。} 交叉验证（cross-validation，CV）在训练数据内部轮流留出一部分，检查模型对没参与本轮拟合的数据表现如何。它帮助选择模型复杂度或正则强度；CV本身不会把一个固定模型直接变好。

\item \textbf{E选。} Bagging用不同的重采样训练集训练多个模型，再投票。若单个模型会随训练样本变化而大幅改变，把这些结果平均能减弱这种波动，尤其适用于容易过拟合的树。

\end{itemize}


这几项以训练和测试的来源、预处理一致为前提。若训练用白天图像、测试全是夜间图像，问题还包含分布变化；此时需要先处理数据差异。有关CV选参后怎样得到最终模型，见\StudyRef{Answers}{mt004}{MT004}。


\textbf{出处：} 2020B期中Quiz Q2，原卷第2页（5分）


\textbf{答案依据：} 2020B期中 Q2 随卷／配套参考答案；条件与解释经复核。
